\begin{table}[h]
\centering
\small
\begin{tabular}{r>{\centering\arraybackslash}p{2.9cm}>{\centering\arraybackslash}p{2.0cm}>{\centering\arraybackslash}p{3.6cm}>{\centering\arraybackslash}p{2.3cm}}
\toprule
Threshold (mm) & accepted as reported & corrected & no qualifying precip in window & no AORC data \\
\midrule
2.5 & 81.0\% & 14.3\% & 2.5\% & 2.1\% \\
5 & 73.9\% & 17.2\% & 6.7\% & 2.1\% \\
\rowcolor{gray!15}7.62\textsuperscript{*} & 66.3\% & 18.3\% & 13.2\% & 2.1\% \\
10 & 59.1\% & 19.4\% & 19.4\% & 2.1\% \\
15 & 46.0\% & 19.3\% & 32.5\% & 2.1\% \\
20 & 36.6\% & 17.8\% & 43.4\% & 2.1\% \\
\bottomrule
\end{tabular}
\caption{\texttt{pluvialCorrectionStatus} share of the released file's 1,149,288 claims, swept over six constant precipitation thresholds. Columns correspond to \texttt{accepted\_as\_reported}, \texttt{corrected}, \texttt{no\_qualifying\_precip\_in\_window}, and, combined, \texttt{no\_aorc\_data\_in\_window} with the threshold-invariant \texttt{before\_aorc\_coverage}. \textsuperscript{*}7.62 mm (shaded) is the released default. Table~\ref{tab:pluvial-correction-status} gives the full $n$ claims breakdown at this default threshold.}
\label{tab:threshold-sweep}
\end{table}
